\documentclass[full-course-notes.tex]{subfiles}

\begin{document}
\lecturemeta{5}
  {Finite Automata: Deterministic and Nondeterministic; Thompson's Construction (Regular
   Expression $\to$ NFA)}
  {CLO 1, CLO 2}
  {\S3.6 (Recognition of Tokens), \S3.7 (Finite Automata; From a Regular Expression to an
   Automaton), Alg.\,3.23 (Thompson's Construction)}
  {[AP] Modern Compiler Implementation, \S2.3--2.4; [HMU] Introduction to Automata Theory,
   Languages, and Computation, Ch.\,2}
\lecshort{Finite Automata \& Thompson's Construction}
\setcounter{section}{0}
\makelecturetitle

\tikzset{
  st/.style={draw, circle, minimum size=0.85cm, font=\small, fill=nuLight, draw=nuNavy},
  fin/.style={st, double, double distance=2pt, fill=dbGreenBg, draw=dbGreen},
  fragst/.style={draw, circle, minimum size=0.7cm, font=\scriptsize, fill=nuLight, draw=nuNavy},
  fragfin/.style={fragst, double, double distance=1.5pt, fill=dbGreenBg, draw=dbGreen},
}

\begin{coreidea}[Lecture Roadmap]
\begin{enumerate}
  \item Why a regular expression alone isn't something a computer can \emph{run} -- and what
  we need instead.
  \item \textbf{Finite automata}, formally: the 5-tuple definition, transition diagrams, and
  what it means for a string to be ``accepted.''
  \item \textbf{DFA vs.\ NFA}: exactly one choice per step vs.\ several (or none) -- and why
  Lecture~\ref{lec:3}'s Chomsky-hierarchy table already told you these have \emph{equal} power.
  \item \textbf{Thompson's construction}: a purely mechanical algorithm that builds an NFA
  directly from the structure of a regular expression, one small piece per operator -- worked
  all the way through on the running example from Lecture~\ref{lec:4}, $(a\mid b)^{*}abb$.
\end{enumerate}
This lecture builds the second arrow in Lecture~\ref{lec:3}'s big-picture diagram: regular
expression $\to$ NFA. The next lecture builds the rest: NFA $\to$ DFA $\to$ minimized DFA
$\to$ an actual running lexer.
\end{coreidea}

\section{Why We Need an Actual Machine}
\label{sec:l5-why-a-machine}

\dbtag A regular expression is a piece of notation -- it describes a language on paper, but a
computer cannot ``run'' a piece of algebra directly. To actually scan a source file character
by character and decide, in real time, which token pattern (if any) the input matches, we need
something a processor can execute: a small, fixed set of states, and a rule for how to move
between them as each character arrives. That is exactly what a \textbf{finite automaton} is,
and Lecture~\ref{lec:3}'s big-picture diagram (\Sref{sec:l3-big-picture}) already named the
first mechanical step that gets us there: \textbf{Thompson's construction}, which converts any
regular expression into an equivalent finite automaton. This lecture builds that step in full;
the next lecture finishes the journey to a working, table-driven lexer.

\section{Finite Automata, Formally}
\label{sec:l5-fa-formal}

\dbtag Before comparing deterministic and nondeterministic versions, it helps to pin down what
a finite automaton actually \emph{is} -- the same formal-definition instinct
\Sref{sec:l3-regular-languages-formal} used for regular languages themselves.

\begin{defbox}{Finite Automaton}
A \textbf{finite automaton} is a 5-tuple $(Q, \Sigma, \delta, q_0, F)$ where:
\begin{itemize}
  \item $Q$ is a finite set of \textbf{states};
  \item $\Sigma$ is the input \textbf{alphabet} (the same notion from
  \Sref{sec:l4-strings-alphabets-languages});
  \item $\delta$ is the \textbf{transition function}, describing how to move between states as
  input symbols arrive (its exact type differs for DFAs and NFAs -- see below);
  \item $q_0 \in Q$ is the \textbf{start state};
  \item $F \subseteq Q$ is the set of \textbf{accepting} (or \emph{final}) states.
\end{itemize}
\end{defbox}

A \textbf{transition diagram} is simply this 5-tuple drawn as a picture: a circle per state, a
double circle for each accepting state, an arrow (labeled with an input symbol) per transition,
and an unlabeled incoming arrow marking the start state -- exactly matching the ``transition
diagram vs.\ transition table'' distinction from \Sref{sec:l3-big-picture}'s exam-focus box:
same underlying automaton, two different ways of writing it down.

\begin{examplebox}{Worked Example: Even Number of 0's}
Recall from \Sref{sec:l4-more-worked-examples} that $1^{*}(01^{*}01^{*})^{*}$ describes binary
strings with an even number of \texttt{0}'s. Here is a two-state automaton for exactly that
language: $Q = \{q_0, q_1\}$, $\Sigma = \{0,1\}$, start state $q_0$, accepting states
$F = \{q_0\}$ (start \emph{and} accept -- zero \texttt{0}'s is even).
\begin{center}
\begin{tikzpicture}[node distance=2.6cm, >=Latex]
  \node[fin] (q0) {$q_0$};
  \node[st, right=of q0] (q1) {$q_1$};
  \draw[trans, ->] (-1.3,0) -- (q0);
  \draw[trans, ->] (q0) to[bend left=25] node[above, font=\small]{0} (q1);
  \draw[trans, ->] (q1) to[bend left=25] node[below, font=\small]{0} (q0);
  \draw[trans, ->] (q0) to[out=135, in=225, looseness=6] node[left, font=\small]{1} (q0);
  \draw[trans, ->] (q1) to[out=45, in=-45, looseness=6] node[right, font=\small]{1} (q1);
\end{tikzpicture}
\end{center}
Reading a \texttt{0} always \emph{flips} which state you're in (tracking odd/even \texttt{0}'s
seen so far); reading a \texttt{1} never changes the count, so it's a self-loop either way. Try
tracing \texttt{1010}: start at $q_0$, read \texttt{1} (stay $q_0$), read \texttt{0} (move
$q_1$), read \texttt{1} (stay $q_1$), read \texttt{0} (move $q_0$) -- ends at $q_0 \in F$,
accepted, matching that \texttt{1010} does indeed have an even number (two) of \texttt{0}'s.
\end{examplebox}

\begin{coreidea}[What ``Accepts'' Means, Formally]
Extend $\delta$ to $\hat\delta$, a function on \emph{strings} rather than single symbols, by
$\hat\delta(q, \varepsilon) = q$ and $\hat\delta(q, xa) = \delta(\hat\delta(q, x), a)$ for a
string $x$ followed by one more symbol $a$ -- in plain words, ``follow the transitions one
symbol at a time, starting from $q$.'' A string $w$ is \textbf{accepted} by the automaton if
$\hat\delta(q_0, w) \in F$: start at the start state, follow $w$ symbol by symbol, and end up
somewhere in the accepting set. The \textbf{language of the automaton}, written $L(\text{automaton})$,
is the set of all strings it accepts -- and for every automaton we build in this course,
$L(\text{automaton})$ will turn out to be exactly the regular language we started with.
\end{coreidea}

\section{DFA vs.\ NFA: Two Models, Equal Power}
\label{sec:l5-dfa-vs-nfa}

\dbtag The automaton above had exactly one arrow leaving each state for each symbol -- no
ambiguity about where to go next. That property has a name.

\begin{defbox}{Deterministic Finite Automaton (DFA)}
A finite automaton is \textbf{deterministic} if $\delta$ is a genuine function
$Q \times \Sigma \to Q$: for every state and every input symbol, there is \emph{exactly one}
transition, and no transitions are labeled with $\varepsilon$. At every step, a DFA's next move
is completely determined -- hence the name.
\end{defbox}

Real regular expressions, though, are built from \emph{union} -- ``match $r$ \emph{or} $s$''
-- which is naturally ambiguous: from a given state, on a given symbol, there might genuinely be
more than one place a matcher could go. Rather than resolve that ambiguity immediately, it is
far more convenient (as Thompson's construction, \Sref{sec:l5-thompson}, will show) to allow it
explicitly in the automaton itself.

\begin{defbox}{Nondeterministic Finite Automaton (NFA)}
A finite automaton is \textbf{nondeterministic} if $\delta$ is a function
$Q \times (\Sigma \cup \{\varepsilon\}) \to 2^{Q}$ (the power set of $Q$): for a given state and
input symbol (or $\varepsilon$, the empty string), there can be \textbf{zero, one, or many}
next states. Two new abilities appear here that a DFA does not have:
\begin{itemize}
  \item \textbf{Multiple choices.} From one state, on one symbol, the automaton may be able to
  move to several different states at once.
  \item \textbf{$\varepsilon$-transitions.} The automaton may move between states
  \emph{without reading any input at all} -- these are the glue Thompson's construction uses to
  wire smaller automata together.
\end{itemize}
A string $w$ is accepted by an NFA if \emph{at least one} sequence of choices, following the
symbols of $w$ (and optionally taking any number of $\varepsilon$-transitions in between), ends
in an accepting state -- not every path has to succeed, just one.
\end{defbox}

\begin{examplebox}{A Small NFA: Strings Ending in \texttt{ab}}
$Q = \{p_0, p_1, p_2\}$, $\Sigma = \{a,b\}$, start $p_0$, accept $F = \{p_2\}$:
\begin{center}
\begin{tikzpicture}[node distance=2.4cm, >=Latex]
  \node[st] (p0) {$p_0$};
  \node[st, right=of p0] (p1) {$p_1$};
  \node[fin, right=of p1] (p2) {$p_2$};
  \draw[trans, ->] (-1.3,0) -- (p0);
  \draw[trans, ->] (p0) to[out=135, in=225, looseness=6] node[left, font=\small]{$a,b$} (p0);
  \draw[trans, ->] (p0) -- node[above, font=\small]{$a$} (p1);
  \draw[trans, ->] (p1) -- node[above, font=\small]{$b$} (p2);
\end{tikzpicture}
\end{center}
From $p_0$, on an \texttt{a}, there are \emph{two} valid moves: stay at $p_0$ (self-loop) or
advance to $p_1$ -- exactly the nondeterminism a DFA cannot express directly. To accept
\texttt{aab}: stay at $p_0$ reading the first \texttt{a} (via the self-loop), then move
$p_0\to p_1$ reading the second \texttt{a}, then $p_1 \to p_2$ reading \texttt{b} -- one
successful path is all it takes, even though \emph{guessing wrong} (always taking the self-loop)
would never reach $p_2$.
\end{examplebox}

\begin{examfocus}
This ``one successful path is enough'' rule is exactly why simulating an NFA directly is
awkward: naively, you'd have to try every possible sequence of choices to know whether
\emph{any} of them accepts. Lecture~\ref{lec:6}'s subset construction sidesteps this entirely,
by tracking \emph{every reachable state at once} as a single set -- turning ``does some path
accept?'' into an ordinary DFA question. That NFA-to-DFA conversion is also the constructive
proof that NFAs and DFAs recognize exactly the same class of languages (Lecture~\ref{lec:3}'s
Chomsky-hierarchy table already told you this fact; Lecture~\ref{lec:6} shows \emph{why} it's
true).
\end{examfocus}

\section{Thompson's Construction: Regular Expression $\to$ NFA}
\label{sec:l5-thompson}

\dbtag Thompson's construction (DB Algorithm 3.23) builds an NFA from a regular expression by
\emph{structural recursion} -- exactly mirroring the six-rule inductive definition of a regular
expression itself from \Sref{sec:l4-regular-expressions}. Every rule below produces a fragment
with exactly \textbf{one} start state and exactly \textbf{one} accept state, which is what
makes the rules compose so cleanly: a bigger fragment is always built by wiring together
smaller fragments' single start/accept points, never by reaching inside them.

\subsection{The Five Construction Rules}
\label{sec:l5-thompson-rules}

\textbf{Base case 1: $\varepsilon$.} A two-state fragment connected by a single
$\varepsilon$-edge.
\begin{center}
\begin{tikzpicture}[node distance=1.8cm, >=Latex]
  \node[fragst] (i) {};
  \node[fragfin, right=of i] (f) {};
  \draw[trans, ->] (i) -- node[above, font=\scriptsize]{$\varepsilon$} (f);
\end{tikzpicture}
\end{center}

\textbf{Base case 2: a single symbol $a \in \Sigma$.} A two-state fragment connected by one
$a$-edge.
\begin{center}
\begin{tikzpicture}[node distance=1.8cm, >=Latex]
  \node[fragst] (i) {};
  \node[fragfin, right=of i] (f) {};
  \draw[trans, ->] (i) -- node[above, font=\scriptsize]{$a$} (f);
\end{tikzpicture}
\end{center}

\textbf{Inductive case: union $r \mid s$.} Given fragments $N(r)$ and $N(s)$ (each with its own
start/accept), build a new start and new accept state, and wire in \emph{both} old fragments
with $\varepsilon$-edges. The old accept states stop being accepting -- only the new one is.
\begin{center}
\begin{tikzpicture}[node distance=1.6cm, >=Latex]
  \node[fragst] (i) {};
  \node[draw, rounded corners, minimum width=1.6cm, minimum height=0.8cm, right=1.3cm of i, yshift=0.9cm, font=\scriptsize, align=center] (Nr) {$N(r)$};
  \node[draw, rounded corners, minimum width=1.6cm, minimum height=0.8cm, right=1.3cm of i, yshift=-0.9cm, font=\scriptsize, align=center] (Ns) {$N(s)$};
  \node[fragfin, right=3.6cm of i] (f) {};
  \draw[trans, ->] (i) -- node[above, sloped, font=\scriptsize]{$\varepsilon$} (Nr.west);
  \draw[trans, ->] (i) -- node[below, sloped, font=\scriptsize]{$\varepsilon$} (Ns.west);
  \draw[trans, ->] (Nr.east) -- node[above, sloped, font=\scriptsize]{$\varepsilon$} (f);
  \draw[trans, ->] (Ns.east) -- node[below, sloped, font=\scriptsize]{$\varepsilon$} (f);
\end{tikzpicture}
\end{center}

\textbf{Inductive case: concatenation $rs$.} Given $N(r)$ and $N(s)$, connect $N(r)$'s accept
state to $N(s)$'s start state with a single $\varepsilon$-edge. $N(r)$'s start becomes the
overall start; $N(s)$'s accept becomes the overall accept.
\begin{center}
\begin{tikzpicture}[node distance=1.6cm, >=Latex]
  \node[draw, rounded corners, minimum width=1.8cm, minimum height=0.8cm, font=\scriptsize, align=center] (Nr) {$N(r)$};
  \node[draw, rounded corners, minimum width=1.8cm, minimum height=0.8cm, right=1.9cm of Nr, font=\scriptsize, align=center] (Ns) {$N(s)$};
  \draw[trans, ->] (Nr.east) -- node[above, font=\scriptsize]{$\varepsilon$} (Ns.west);
\end{tikzpicture}
\end{center}

\textbf{Inductive case: closure $r^{*}$.} Given $N(r)$, build a new start and new accept state.
Wire the new start directly to the new accept ($\varepsilon$, for zero repetitions) and to
$N(r)$'s start ($\varepsilon$, to begin a repetition); wire $N(r)$'s accept back to $N(r)$'s
start ($\varepsilon$, to loop again) and forward to the new accept ($\varepsilon$, to stop).
\begin{center}
\begin{tikzpicture}[node distance=1.6cm, >=Latex]
  \node[fragst] (i) {};
  \node[draw, rounded corners, minimum width=1.8cm, minimum height=0.8cm, right=1.6cm of i, font=\scriptsize, align=center] (Nr) {$N(r)$};
  \node[fragfin, right=1.6cm of Nr] (f) {};
  \draw[trans, ->] (i) -- node[above, font=\scriptsize]{$\varepsilon$} (Nr.west);
  \draw[trans, ->] (Nr.east) -- node[above, font=\scriptsize]{$\varepsilon$} (f);
  \draw[trans, ->] (i) to[bend right=45] node[below, font=\scriptsize]{$\varepsilon$} (f);
  \draw[trans, ->] (Nr.south) to[out=250, in=290, looseness=6] node[below, font=\scriptsize]{$\varepsilon$} (Nr.south west);
\end{tikzpicture}
\end{center}

\begin{examfocus}
Every fragment above has exactly \textbf{one} start and \textbf{one} accept state -- even after
combining fragments. This invariant is what makes the recursion work at all: at every step, you
always know exactly which single state to wire the \emph{next} $\varepsilon$-edge into, no
matter how deeply nested the regular expression is. Losing track of this (e.g.\ trying to wire
a new edge into ``one of the accept states'' of a union fragment) is the most common mistake
when doing this construction by hand.
\end{examfocus}

\subsection{Worked Example: Building the NFA for $(a\mid b)^{*}abb$}
\label{sec:l5-thompson-worked}

\dbtag This is exactly the regular expression \Sref{sec:l4-worked-example-build-up} built up
piece by piece with a syntax tree. We now build its NFA the same way, stage by stage, applying
one construction rule per stage.

\begin{examplebox}{Stage 1 -- The Atoms}
Build the two-state fragments for the individual symbols first. We'll need four of them: one
for \texttt{a} inside the $(a\mid b)$, one for \texttt{b} inside the $(a\mid b)$, and one more
each for the trailing \texttt{a}, \texttt{b}, \texttt{b}.
\begin{center}
\begin{tikzpicture}[node distance=1.6cm, >=Latex]
  \node[fragst] (1) {1};
  \node[fragfin, right=of 1] (2) {2};
  \draw[trans, ->] (1) -- node[above, font=\scriptsize]{$a$} (2);
  \node[fragst, right=1.6cm of 2] (3) {3};
  \node[fragfin, right=of 3] (4) {4};
  \draw[trans, ->] (3) -- node[above, font=\scriptsize]{$b$} (4);
\end{tikzpicture}
\end{center}
(States 1--2 recognize just \texttt{a}; states 3--4 recognize just \texttt{b}.)
\end{examplebox}

\begin{examplebox}{Stage 2 -- Union: $(a\mid b)$}
Apply the union rule to the two fragments above: new start state 5, new accept state 6.
\begin{center}
\begin{tikzpicture}[node distance=1.5cm, >=Latex]
  \node[fragst] (5) {5};
  \node[fragst, right=2.3cm of 5, yshift=0.9cm] (1) {1};
  \node[fragfin, right=of 1] (2) {2};
  \node[fragst, right=2.3cm of 5, yshift=-0.9cm] (3) {3};
  \node[fragfin, right=of 3] (4) {4};
  \node[fragfin, right=2.3cm of 2] (6) {6};
  \draw[trans, ->] (5) -- node[above, sloped, font=\scriptsize]{$\varepsilon$} (1);
  \draw[trans, ->] (5) -- node[below, sloped, font=\scriptsize]{$\varepsilon$} (3);
  \draw[trans, ->] (1) -- node[above, font=\scriptsize]{$a$} (2);
  \draw[trans, ->] (3) -- node[above, font=\scriptsize]{$b$} (4);
  \draw[trans, ->] (2) -- node[above, sloped, font=\scriptsize]{$\varepsilon$} (6);
  \draw[trans, ->] (4) -- node[below, sloped, font=\scriptsize]{$\varepsilon$} (6);
\end{tikzpicture}
\end{center}
(States 2 and 4 are no longer accepting -- only the new state 6 is.)
\end{examplebox}

\begin{examplebox}{Stage 3 -- Closure: $(a\mid b)^{*}$}
Apply the closure rule to the union fragment above (states 5--6): new start state 7, new accept
state 8.
\begin{center}
\begin{tikzpicture}[node distance=1.4cm, >=Latex]
  \node[fragst] (7) {7};
  \node[draw, rounded corners, minimum width=2.6cm, minimum height=1.9cm, right=1.6cm of 7, font=\scriptsize, align=center] (U) {$(a\mid b)$\\\emph{states 5--6}};
  \node[fragfin, right=1.6cm of U] (8) {8};
  \draw[trans, ->] (7) -- node[above, font=\scriptsize]{$\varepsilon$} (U.west);
  \draw[trans, ->] (U.east) -- node[above, font=\scriptsize]{$\varepsilon$} (8);
  \draw[trans, ->] (7) to[bend right=50] node[below, font=\scriptsize]{$\varepsilon$ (skip)} (8);
  \draw[trans, ->] (U.south) to[out=260, in=280, looseness=5] node[below, font=\scriptsize]{$\varepsilon$ (loop)} (U.south west);
\end{tikzpicture}
\end{center}
State 7 is now the start of everything built so far; state 8 is the accept state for
$(a\mid b)^{*}$.
\end{examplebox}

\begin{examplebox}{Stages 4--6 -- Concatenating $a$, $b$, $b$}
Three more single-symbol fragments (states 9--10 for \texttt{a}, 11--12 for \texttt{b}, 13--14
for \texttt{b}), concatenated on one at a time: state 8 connects to state 9, state 10 connects
to state 11, state 12 connects to state 13 -- each via a single $\varepsilon$-edge, exactly the
concatenation rule.
\begin{center}
\begin{tikzpicture}[node distance=1.3cm, >=Latex]
  \node[draw, rounded corners, minimum width=2.6cm, minimum height=1.1cm, font=\scriptsize, align=center] (star) {$(a\mid b)^{*}$\\\emph{states 7--8}};
  \node[fragst, right=1.4cm of star] (9) {9};
  \node[fragfin, right=of 9] (10) {10};
  \node[fragst, right=1.4cm of 10] (11) {11};
  \node[fragfin, right=of 11] (12) {12};
  \node[fragst, right=1.4cm of 12] (13) {13};
  \node[fragfin, right=of 13] (14) {14};
  \draw[trans, ->] (star.east) -- node[above, font=\scriptsize]{$\varepsilon$} (9);
  \draw[trans, ->] (9) -- node[above, font=\scriptsize]{$a$} (10);
  \draw[trans, ->] (10) -- node[above, font=\scriptsize]{$\varepsilon$} (11);
  \draw[trans, ->] (11) -- node[above, font=\scriptsize]{$b$} (12);
  \draw[trans, ->] (12) -- node[above, font=\scriptsize]{$\varepsilon$} (13);
  \draw[trans, ->] (13) -- node[above, font=\scriptsize]{$b$} (14);
\end{tikzpicture}
\end{center}
The complete NFA for $(a\mid b)^{*}abb$ has 14 states, start state 7, and a single accept
state, 14. Every edge in every stage above is still present -- concatenation only ever
\emph{adds} an edge, it never removes or renames anything from the fragments it connects.
\end{examplebox}

\begin{examplebox}{Tracing \texttt{"babb"} Through the Finished NFA}
Recall from \Sref{sec:l4-worked-example-build-up} that \texttt{babb} should be accepted. One
successful path: $7 \xrightarrow{\varepsilon} 5 \xrightarrow{\varepsilon} 3 \xrightarrow{b} 4
\xrightarrow{\varepsilon} 6 \xrightarrow{\varepsilon} 8 \xrightarrow{\varepsilon} 9
\xrightarrow{a} 10 \xrightarrow{\varepsilon} 11 \xrightarrow{b} 12 \xrightarrow{\varepsilon} 13
\xrightarrow{b} 14$. Reading off just the symbols consumed along this path --
$b, a, b, b$ -- reproduces \texttt{babb} exactly, and the path ends at state 14, the accept
state. One successful path is all an NFA needs.
\end{examplebox}

\section{Self-Check Questions}
\label{sec:l5-selfcheck}

\begin{enrichment}[Practice -- Not Graded]
\begin{enumerate}
  \item Write out the formal 5-tuple $(Q, \Sigma, \delta, q_0, F)$ for the ``even number of
  0's'' automaton in \Sref{sec:l5-fa-formal}, giving $\delta$ as an explicit table (one row per
  state, one column per symbol).
  \item Using $\hat\delta$ from the coreidea box in \Sref{sec:l5-fa-formal}, trace whether the
  string \texttt{0110} is accepted by the same automaton. Show each intermediate state.
  \item Explain, without using the word ``choice,'' what it means for an NFA's transition
  function to have type $Q \times (\Sigma \cup \{\varepsilon\}) \to 2^{Q}$ rather than
  $Q \times \Sigma \to Q$.
  \item In the ``strings ending in \texttt{ab}'' NFA (\Sref{sec:l5-dfa-vs-nfa}), list \emph{all}
  paths (successful or not) that consume the string \texttt{aab}, and confirm exactly one of
  them ends in the accepting state.
  \item Apply Thompson's construction by hand to the regular expression $ab^{*}$: draw the
  fragment for \texttt{a}, the fragment for $b^{*}$, and the concatenation that joins them.
  \item Why does every fragment in Thompson's construction have \emph{exactly one} start state
  and \emph{exactly one} accept state, even after several rounds of union, concatenation, and
  closure? What would go wrong with the concatenation rule if a fragment were allowed more than
  one accept state?
  \item For the finished 14-state NFA in \Sref{sec:l5-thompson-worked}, find one string of
  length 5 that is accepted and one string of length 5 that is \emph{not} accepted, and briefly
  justify each.
\end{enumerate}
\end{enrichment}

\begin{takeaways}
\begin{itemize}
  \item A \textbf{finite automaton} is a 5-tuple $(Q, \Sigma, \delta, q_0, F)$; a
  \textbf{transition diagram} is just its picture, and a \textbf{transition table} is the same
  information as an array -- interchangeable representations, not different machines.
  \item A \textbf{DFA} has exactly one transition per state per symbol and no
  $\varepsilon$-edges; an \textbf{NFA} allows several transitions (or none) per state per
  symbol, plus $\varepsilon$-edges that move without consuming input.
  \item An NFA accepts a string if \emph{at least one} sequence of choices reaches an accepting
  state -- not every path has to succeed.
  \item \textbf{Thompson's construction} builds an NFA from a regular expression using five
  rules ($\varepsilon$, a single symbol, union, concatenation, closure), each producing a
  fragment with exactly one start and one accept state -- which is exactly what lets the rules
  compose recursively, mirroring the regular expression's own inductive definition rule for
  rule.
  \item The worked $(a\mid b)^{*}abb$ NFA (14 states) is the same running example from
  Lecture~\ref{lec:4} -- Lecture~\ref{lec:6} converts this exact automaton into a DFA and then
  minimizes it, so it is worth understanding this construction solidly before moving on.
\end{itemize}
\end{takeaways}

\section*{References for This Lecture}
\begin{thebibliography}{99}
\bibitem[DB]{db} Aho, A.~V., Lam, M.~S., Sethi, R., and Ullman, J.~D. \textit{Compilers:
Principles, Techniques, and Tools}. 2nd~ed. Pearson, 2006.
\bibitem[AP]{ap} Appel, A.~W. \textit{Modern Compiler Implementation in C/Java/ML}. Cambridge
University Press, 2004.
\bibitem[HMU]{hmu} Hopcroft, J.~E., Motwani, R., and Ullman, J.~D. \textit{Introduction to
Automata Theory, Languages, and Computation}. 3rd~ed. Pearson, 2006.
\end{thebibliography}

\end{document}
